% MAS1004 Assignment 2: long report.
%
% You fill this in with your agent, in your repository. AGENTS.md tells the
% agent how. Red text in the PDF is what still has to be replaced.
%
% - Pictures go in figures/. Copy them from results/, which is not in git.
% - Outputs go in outputs/: copy experiments.csv, LOG.md and the outputs of
%   check_submission.py and curate.py count here. \outputfile{outputs/...}
%   prints a file exactly as it is. Every number in this report has to come
%   from such an output.
%
% To make the PDF: compress this folder into a zip, and on Overleaf choose
% New Project, Upload Project. Set Menu, Compiler to XeLaTeX, and Recompile.

\documentclass[11pt]{article}
\usepackage{mas1004}

\reporttitle{Assignment 2: Long report}
\studentname{\fillin{your name}}
\studentid{\fillin{your student ID}}

\begin{document}
\makeheader

Leaderboard nickname \textbf{Demogom}. My last releases are \texttt{track-a-2}
(moving average of the weights) and \texttt{track-b-1} (\texttt{train\_b.py} as
given).

% ---------------------------------------------------------------------------
\section*{Track A}

\begin{center}\footnotesize
\begin{tabular}{lp{0.6\linewidth}r}
\toprule
tag & what changed from the run before & validation \\
\midrule
baseline & the starter recipe, nothing changed & 0.4987 \\
baseline, seed 1 & nothing at all: the same recipe with the other seed, to
see how much the score moves by chance & 0.4971 \\
busan20 & every Busan photo was drawn 1.88 times as often, so Busan gets
20.0\% of the 250,000 photos instead of 11.7\%; the other four cities keep
their proportions & 0.4889 \\
noflip & \texttt{T.RandomHorizontalFlip()} deleted, Busan weighting back to
the starter recipe & 0.4889 \\
ema & the starter recipe again, but the model that is handed in is a moving
average of the weights over the whole run (decay 0.999) instead of the
weights of the last step & 0.5052 \\
\bottomrule
\end{tabular}
\end{center}

I ran the starter recipe twice before changing anything. The two seeds gave
0.4987 and 0.4971, so the score moves by 0.0016 when nothing changes at all.
I took that as my measuring stick: a change smaller than about 0.002 is luck,
not a result, and I read every later run against it.

\textbf{busan20.} Busan has 5,191 training photos against Tokyo's 13,347, so
the starter recipe puts 11.7\% of the budget on Busan and 30.2\% on Tokyo. My
reason for the change was that Busan having less than half of Tokyo's photos
seemed unfair to Busan. With a fixed budget of 250,000 photos this is zero
sum, so the run is also a test of whether the weakest city gains more than the
others lose. It did not: Busan went from 0.378 / 0.402 to \textbf{0.4374}, but
Seoul fell from 0.600 to 0.5441, Osaka from 0.351 to 0.3365 and Taipei from
0.535 to 0.5033, and the total fell by 0.009, five times the seed noise. My
reading: the Busan photos pushed the other four cities out of the budget, so
the model stopped paying attention to them.

\textbf{noflip.} Tokyo and Osaka drive on the left, Seoul, Busan and Taipei on
the right, so a horizontally flipped photo shows lanes, signs and traffic that
cannot exist in that city; I expected removing the flip to help the three
right-hand cities and Osaka. It gave \textbf{0.4889}, below both seeds. Only
Taipei moved up (0.527 against 0.535 / 0.509); Seoul, Busan and Osaka all fell.
So the flip was not confusing the model, and it was probably doing useful
regularisation work.

\textbf{ema.} The gap between my validation score (0.4979) and my leaderboard
score (0.409) is about 0.09, so the model fits the starter photos much better
than it fits frames from channels it has never seen. A cheap way to attack
that, without touching the photos, is to hand in a moving average of the
weights: the one-cycle schedule ends with a nearly zero learning rate, so the
last weights are close to the average of the end of training already, and the
average removes what is left of the noise. It gave \textbf{0.5052}, 0.007
above the mean of the two seeds and past 0.5. Seoul (0.6074) and Tokyo (0.6478)
are the two cities that moved, Osaka (0.3689) a little, Busan (0.382) and
Taipei (0.5198) slightly down.

\textbf{My final recipe and selection.} The recipe is the starter recipe with
the moving average (\texttt{ema}); the selection is all 44,257 starter training
photos, so I have no \texttt{selection.csv} and no photo was dropped. Its
validation score is 0.5052. On the leaderboard the starter recipe (my
\texttt{track-a-1}) scored \textbf{0.409} while validation said 0.4987;
\texttt{track-a-2} was released the night before this report was written, and
the grader takes the newest release of each day, so its leaderboard number
arrives afterwards. The two numbers disagree for a reason I can name: my
validation photos come from the same places as my training photos (Wikimedia
Commons, walking and driving videos of the same kind), the leaderboard uses
frames from channels that are not in the starter data at all, and the final
ranking uses phone street photos. All three are street scenes, but only the
first is the data the model was fitted to.

% ---------------------------------------------------------------------------
\section*{Decisions you made from photos}

\picturefile{figures/mistakes_baseline_seed0_busan_00.png}

\picturefile{figures/mistakes_baseline_seed0_osaka_00.png}

I looked at the photos the baseline model got wrong with the most confidence,
48 per city (\texttt{curate.py mistakes}). The counts were seoul 974 of 2,435,
busan 864 of 1,390, tokyo 1,440 of 3,881, osaka 1,382 of 2,128 and taipei 1,042
of 2,243, and the model answers Seoul very often when it is wrong. Two things
the wrong photos had in common, in my words: streets with a lot of people, and
videos taken from a car. The rule I wrote from that was to be careful with
crowds and with road-view frames.

\fillin{what I decided to do with that rule, what it changed and the scores
before and after -- to fill in with the student}

\fillin{the requests I gave my agent for these decisions}

% ---------------------------------------------------------------------------
\section*{Track B}

I started with \texttt{train\_b.py} as it is given: ResNet18 at 224 pixels,
4 epochs, batch 64, learning rate 3e-4, on all 44,257 starter training photos,
which is 176,896 photos through the network, and it took 142 minutes on this
computer's CPU. The validation score was \textbf{0.4952} (seoul 0.5708,
busan 0.3755, tokyo 0.6189, osaka 0.344, taipei 0.5667), close to my Track A
starter recipe, which sees 250,000 photos for 0.4979. The model is 11.17
million parameters and 1.81 billion multiply-adds, so it fits the limits.

I collected no extra photos yet; the next step I would take is phone photos
sorted into the city folders with \texttt{gps\_sort.py}.

On the leaderboard this release is too new to score in time for this report:
\texttt{track-b-1} was released the night before. The Track A comparison is
the same shape as here: validation 0.4987 became 0.409 on the leaderboard,
because the leaderboard frames come from channels that are not in the starter
data and the final ranking uses phone photos.

\subsection*{Your submission, as the grader sees it}
\outputfile{outputs/check_submission.txt}

% ---------------------------------------------------------------------------
\section*{What made the difference}

The only change that helped was moving the weights to their own average: +0.007
on validation, 3.5 times the seed noise, with the photos and the schedule left
exactly as they were. It is also the cheapest change I made, since it costs no
training time. Removing the horizontal flip did not help at all (0.4889): the
flip was not confusing the model. Re-weighting the cities hurt most (0.4889):
it lifted the city I aimed at by 0.047 and pulled the other four down by more,
because a fixed budget of 250,000 photos cannot give one city more without
taking from another.

With one more week I would:

\begin{itemize}
\item collect real phone and street photos, especially for Busan (it is the
  only city with no camera-on-a-car photos at all) and Taipei;
\item act on the rule from the mistakes sheets and drop the non-street photos,
  and measure it, since that is the one lever I looked at but did not try;
\item run the \texttt{ema} recipe twice more with other seeds, to be sure
  +0.007 is bigger than luck;
\item try mixup on top of \texttt{ema} in Track A;
\item in Track B, MobileNetV3-Large at 320 pixels with balanced sampling (4.2
  million parameters, 0.44 billion multiply-adds, so it fits), which is the
  kind of change the limits leave room for.
\end{itemize}

% ---------------------------------------------------------------------------
\section*{All runs}
\outputfile{outputs/experiments.csv}

\section*{Lab notebook}
\outputfile{outputs/LOG.md}

\end{document}
